% Shared exercise chunk: l1-greedy-decoding
\begin{exercisechunk}
\item \textbf{Greedy decoding and repetition.}
\label[exercise]{ex:greedy-decoding}
A toy language model completes the prefix ``The movie was''.
Let ``\texttt{very}'' and ``\texttt{good.}'' (including the full stop, not including the apostrophes, which we will drop)
be two tokens in the model's output token vocabulary.
Given this prefix, the toy model outputs
token \texttt{very} with probability $0.6$ or
token \texttt{good.}
with probability $0.4$.
When the prefix ends with token
\texttt{very}, it repeats the same
choice; when it ends with \texttt{good.},
it outputs the stop token with probability one and generation stops.
Count the stop token as a generated token, including for the budget below.

\begin{enumerate}[(a)]
\item What token sequence does greedy decoding produce with a budget of $10$
additional tokens? What token sequence does it produce without a budget?
\item Under sampling, calculate the probability that generation terminates
within $10$ additional tokens. Without a token budget, show that generation
terminates with probability one and compute the expected number of generated
tokens.
\item Which complete sentence has the highest probability under the model?
What is its probability?
\end{enumerate}

\paragraph{Connection to language models.}
\href{https://arxiv.org/pdf/1908.04319v2}{Welleck et al. (2019, Table~4,
p.~11)} show a GPT-2 continuation that repeats ``Portuguese'' under greedy
decoding. The
\href{https://huggingface.co/Qwen/Qwen3-8B}{Qwen3-8B model card}
also warns that greedy decoding can cause endless repetition in thinking mode.
\end{exercisechunk}
